% Pursuing Stacks (cotes 134-2 to 134-8): archive cote 134-2, site page 16 (PDF page 17).
% letter to Quillen. Draft compared: ps1.tex:155-228.
% AI-transcribed by Claude Opus 5.5 (claude-opus-5-5, adaptive thinking, effort high), Message Batches API, batch ps1 (msgbatch_01Rv5hCaNTtUt4JHx1yqoHPL), 2026-11-12; not yet checked by eye.
\scanpage{134-2/16}
\mspage{\hand{4}}

and for $\hand{\ell}=1$
\[ s_1(v \underset{1}{*} u ) = s_1(u) \]
\[ t_1(v \underset{1}{*} u ) = t_1(v) \]
(NB the operation $v \underset{1}{*} u$ is \inserted{just} the usual composition $v \circ u$ ).

One may be tempted to think that the preceding data exhaust the
structure of $\infty$-groupoid, and that they will have to be supplemented only
by a handful of suitable \uline{axioms}, one being \uline{associativity} for the opera-%
tion $\overset{i}{\underset{\hand{\ell}}{*}}$ , which can be expressed essentially by saying that \unsure{that|this} composi
\hand{-}tion operation turns $F_i$ into the set of arrows of a category having $F_{i-\hand{\ell}}$
as set of objects (with\xout{e} the source and target maps $s^i_{\hand{\ell}}$ and $t^i_{\hand{\ell}}$, and
with identity map $k^{i-\hand{\ell}}_{\hand{\ell}} : F_{i-\hand{\ell}} \to F_i$ the composition of the identity
maps $F_{i-\hand{\ell}} \to F_{i-\hand{\ell}+1} \to \ldots F_{i-1} \to F_i$ ), and another being the Godement
relation
\[ ( v' \underset{\hand{\alpha}}{*} v ) \underset{\hand{\nu}}{*} (u' \underset{\hand{\alpha}}{*} u ) = (v' \underset{\hand{\nu}}{*} u' ) \underset{\hand{\alpha}}{*} (v \underset{\hand{\nu}}{*} u ) \]
(with the assumptions \hand{$1\leq\alpha\leq\nu$,} \unsure{and} \hand{$u, u', v, v' \in F_i$ and}
\[ \hand{$\left\{\begin{aligned} t_\alpha(u)&=s_\alpha(u') \\ t_\alpha(v)&=s_\alpha(v') \end{aligned}\right. \qquad t_\nu(u)=s_\nu(v)=s_\nu(\unsure{u'|v'})=t_\nu(u')$} \]
\unsure{\hand{,}}
\deleted{implying} \inserted{expressing} that both members are defin\xout{\unsure{i}}ed ) , plus the two relations
concerning the invers\xout{\unsure{a}}ion of i-objects ($i\geq 1$) $u \mapsto \overset{\hand{\vee}}{u}$ ,
\[ u \underset{1}{*} \overset{\hand{\vee}}{u} = \deleted{\xout{xxxxxx}_{1}} \mathrm{id}_{t_1(u)} \ , \quad \overset{\hand{\vee}}{u} \underset{1}{*} u = \mathrm{id}_{s_1(u)} \]
\[ (\overset{\hand{\vee}}{v} \underset{\hand{\ell}}{*} \overset{\hand{\vee}}{u}) = \overset{\hand{\vee}}{\unsure{\over{u}{v}}} \underset{\hand{\ell}}{*} \overset{\hand{\vee}}{u} \quad (\hand{\ell} \geq 2) \ . \]
It just occurs to me, by the way, that the \inserted[typed]{previous} description of basic (or
``primary'') data for an $\infty$- groupoid is already uncomplete in some rather
obvious respect, namely that the symmetry-operation \inserted{$\mathrm{inv}_i :$} $u \mapsto \overset{\hand{\vee}}{u}$ on $F_i$ must
be comple\xout{\unsure{h}}mented by i-1 similar involutions on $F_i$ , which corresponds
algebraically to the intuition that when we have an \inserted{$(i+1)$-} arrow \hand{$\lambda$} say between
two \inserted{$i$-} arrows $u$ and $v$, then we must be able to deduce from it another arrow
from $\overset{\hand{\vee}}{u}$ to $\overset{\hand{\vee}}{v}$ (namely $u \mapsto \overset{\hand{\vee}}{u}$ has a``functorial character''\xout{ixxx} for variable
$u$ )? This seems a rather anodine modification of the previous set-up,
and is irrelevant for the main point I want to make here, namely: that
for the notion of $\infty$ -groupoids we are after, all the equalities just
envisioned in this paragraph (and those I guess which will ensu\unsure{e} natural-%
ly by the necessary extension of the basic involution on $F_i$ ) should be
replaced by ``homotopies''\hand{,} namely by $(i+1)$-arrows between the two members.
These arrows should be viewed, I believe, as being part of the data, \xout{thx}
\xout{Thexdiffi} they appear here as a kind of ``secondary'' structure. The diffi-%
culty which appears now is to work out the natural coherence properties

\mspage{\hand{16}}
